\section{条件有偏估计：线性收缩与经验贝叶斯}
\label{sec:part13}

\subsection{突破条件无偏下确界的理论可能性}

在定理 4.1 中，下确界 $v_R$ 的刚性约束建立在\textbf{条件无偏性}（$\E(\tilde{\theta} \mid R_t) = R_t$）这一前提假设之上。在现代统计学与机器推断中，诸如 James--Stein 收缩估计量、经验贝叶斯（Empirical Bayes）以及状态空间平滑算法（如 EpiFilter，\citet{parag2021improved}）允许引入确定性条件偏差，以期换取方差的大幅缩减。

\subsection{线性收缩估计量的均方误差闭式推导}

\begin{derivationbox}{推导 13.1：最优收缩强度与先验偏差容限}
构造经典凸组合线性收缩估计量：
\begin{equation}
  \hat{\theta}_\lambda \equiv (1 - \lambda) \frac{C_t}{m} + \lambda \mu_0, \qquad \lambda \in [0, 1],
  \label{eq:shrink-def}
\end{equation}
其中 $\mu_0$ 为先验锚定中心（如多季历史基准或流行病学常模），$\lambda$ 为收缩强度。记先验中心与真实环境均值的偏差为 $\delta \equiv \mu_0 - \Rbar$，无偏估计量方差为 $V \equiv v_R + A/m$。
计算总均方误差：
\begin{align}
  \hat{\theta}_\lambda - \Rbar &= (1 - \lambda)\left( \frac{C_t}{m} - \Rbar \right) + \lambda(\mu_0 - \Rbar) \nonumber \\
  &= (1 - \lambda)\left( \frac{C_t}{m} - \Rbar \right) + \lambda \delta.
\end{align}
由于 $\E(C_t/m - \Rbar) = 0$，平方求期望时交叉项消失：
\begin{equation}
  \text{MSE}(\lambda) \equiv \E\big[(\hat{\theta}_\lambda - \Rbar)^2\big] = (1 - \lambda)^2 V + \lambda^2 \delta^2.
\end{equation}
对收缩参数 $\lambda$ 求一阶导数寻找极值：
\begin{equation}
  \frac{\mathrm{d}\,\text{MSE}(\lambda)}{\mathrm{d}\lambda} = -2(1 - \lambda)V + 2\lambda \delta^2 = 0 \implies \lambda(V + \delta^2) = V.
\end{equation}
由此得出最优收缩权重（Oracle shrinkage intensity）：
\begin{equation}
  \boxed{\, \lambda^* = \frac{V}{V + \delta^2} \,}
  \label{eq:opt-lambda}
\end{equation}
将最优权重代回均方误差公式：
\begin{align}
  \text{MSE}(\lambda^*) &= \left( 1 - \frac{V}{V+\delta^2} \right)^2 V + \left( \frac{V}{V+\delta^2} \right)^2 \delta^2 \nonumber \\
  &= \left( \frac{\delta^2}{V+\delta^2} \right)^2 V + \frac{V^2 \delta^2}{(V+\delta^2)^2} \nonumber \\
  &= \frac{V \delta^4 + V^2 \delta^2}{(V+\delta^2)^2} = \frac{V \delta^2 (V + \delta^2)}{(V+\delta^2)^2} = \frac{V \delta^2}{V + \delta^2}.
  \label{eq:opt-mse-shrink}
\end{align}
\end{derivationbox}

\begin{theorembox}{定理 13.1（收缩估计击穿环境方差的充要条件）}
最优收缩估计量的均方误差严格小于不可约环境方差 $v_R$（即 $\text{MSE}(\lambda^*) < v_R$）当且仅当先验偏差满足：
\begin{equation}
  \delta^2 < v_R \left( 1 + \frac{m}{\mx} \right).
  \label{eq:shrink-barrier}
\end{equation}
\end{theorembox}

\begin{proof}
根据式~\eqref{eq:opt-mse-shrink}，求解不等式：
\begin{equation}
  \frac{V \delta^2}{V + \delta^2} < v_R \iff V \delta^2 < v_R V + v_R \delta^2 \iff (V - v_R) \delta^2 < v_R V.
\end{equation}
已知 $V - v_R = A/m$，且 $V = v_R + A/m$。代入得：
\begin{equation}
  \frac{A}{m} \delta^2 < v_R \left( v_R + \frac{A}{m} \right).
\end{equation}
两侧同乘以 $\frac{m}{A}$：
\begin{equation}
  \delta^2 < v_R \left( \frac{v_R m}{A} + 1 \right).
\end{equation}
利用方差平衡规模定义 $\mx = A / v_R \implies \frac{v_R}{A} = \frac{1}{\mx}$，立即得出：
\begin{equation}
  \delta^2 < v_R \left( 1 + \frac{m}{\mx} \right).
\end{equation}
定理获证。
\end{proof}

\begin{insightbox}{权衡与警示：收缩方差红利的代价}
式~\eqref{eq:shrink-barrier} 证明：只要先验中心 $\mu_0$ 与总体均值的偏差平方 $\delta^2$ 落在容许椭球内，经验贝叶斯确实可以战胜环境方差 $v_R$。
在基准参数（$m = 50, v_R = 0.04$）下，只要先验偏差 $|\delta| < 0.279$，收缩估计的 MSE 就低于 $0.04$；当偏差仅为 $0.1$ 时，MSE 甚至只有 $0.0089$（仅为 $v_R$ 的 22\%）。

\textbf{但天上没有免费的午餐。其惨痛代价是引入了状态依赖的条件偏差：}
\begin{equation}
  \text{Bias}(\hat{\theta}_\lambda \mid R_t) \equiv \E(\hat{\theta}_\lambda \mid R_t) - R_t = (1 - \lambda) R_t + \lambda \mu_0 - R_t = \lambda(\mu_0 - R_t).
\end{equation}
条件偏差随当期传播环境 $R_t$ 偏离常模的幅度\textbf{线性暴增}！在常态稳定期（$R_t \approx \Rbar$），估计量非常精准；而在新发突变株引发超级暴发或断崖式骤降的非常态危机时刻（$|R_t - \mu_0| \gg 0$），收缩估计量因强烈被锚定在历史中心，发生灾难性的钝化与误判。这恰恰违背了公共卫生危机干预的初衷。
\end{insightbox}
